\clearpage
\section{Verification and Testing}

The project combines JUnit tests on the production code with Java PathFinder
(JPF) exploration of concurrent synchronization protocols.

\subsection{JUnit tests}

The Maven suite checks deterministic physics, collisions, pocketing, FIFO
commands, concurrent submissions, failures, snapshots, shutdown, and view
input. JPF uses a Java 11-compatible runtime or the dedicated Docker workflow,
so it is executed separately from the default Java 17 Maven build.

\subsection{JPF verification}

JPF exhaustively explores the schedules of four bounded validation harnesses:

\begin{itemize}
    \item \textbf{\texttt{ThreadedMiniHarness}:} a reduced game-loop model.
    It checks exclusive logical board ownership, exact-once command execution,
    and snapshot publication only after command draining and two worker
    completions.
    \item \textbf{\texttt{TaskBasedMiniHarness}:} the same reduced game-loop
    protocol, with two short-lived task threads instead of long-lived workers.
    \item \textbf{\texttt{PhysicsWorkersJpfHarness}:} directly exercises the
    production \texttt{PhysicsWorker} and \texttt{WorkerCompletionMonitor}.
    It checks that both workers complete before commit, are reused for a second
    phase, and then terminate.
    \item \textbf{\texttt{TaskBasedPhysicsBatchHarness}:} a validation-only
    model of the task-based physics phase. It checks submit-all, wait-all, and
    commit-after-completion for two tasks, including pool reuse and shutdown.
\end{itemize}

\begin{center}
\small
\begin{tabular}{|>{\bfseries}p{3.2cm}|r|r|}
    \hline
    Model & New states & Max. depth \\
    \hline
    Threaded loop & 173\,825 & 206 \\
    \hline
    Task-based loop & 49\,377 & 155 \\
    \hline
    Threaded physics & 13\,153 & 144 \\
    \hline
    Task-based physics & 20\,945 & 103 \\
    \hline
\end{tabular}
\end{center}

All four searches completed with \texttt{no errors detected}: no assertion
failure or deadlock was found in the represented finite scenarios. The result
does not prove numerical physics, full \texttt{ExecutorService} behavior, or
arbitrary workloads; these are covered by production tests and benchmarks.
